% Source-aligned English for printed p.155.
\SourcePage{155}

\LayerTarget{tohoku:construction:sheafification:p2}
\noindent \textit{functor} of $F$, and the homomorphism
$F\to\mathfrak F(\overline F)$ is natural). To interpret a sheaf
of groups, abelian groups, and so on in terms of étalé spaces, one equips
every fiber of the étalé space with a group law, respectively an abelian
group law and so on, satisfying a natural continuity condition. This recovers
the definition of [4, XIV].
\LayerTarget{tohoku:remark:stalkwise-exactness}
--- The corresponding étalé spaces immediately show when a homomorphism of
sheaves $F\to G$ is a monomorphism or an epimorphism: it is
necessary and sufficient that on every fiber $F(x)$ the corresponding
homomorphism $F(x)\to G(x)$ be a monomorphism or an epimorphism,
respectively. Likewise, the fiber at $x$ of the kernel, cokernel, image, or
coimage of a homomorphism of sheaves of abelian groups is obtained by taking
the corresponding kernel, cokernel, and so on of
$F(x)\to G(x)$.

\LayerTarget{tohoku:def:sheaf-of-modules}
Let $\mathbf O$ be a sheaf of rings with identity on $X$. A
\textit{sheaf of left $\mathbf O$-modules}, or briefly a \textit{left
$\mathbf O$-module} on $X$, is a sheaf of abelian groups $F$ together with,
for every open $U\subset X$, a unital left $\mathbf O(U)$-module structure on
$F(U)$ such that the restriction maps $F(U)\to F(V)$ are
compatible with the actions of $\mathbf O(U)$ and $\mathbf O(V)$ on $F(U)$
and $F(V)$, respectively. Equivalently, one gives a homomorphism from the
sheaf of rings $\mathbf O$ to the sheaf of rings of germs of endomorphisms of
the sheaf of abelian groups $F$.
Right $\mathbf O$-modules are defined analogously. For short, we shall simply
say $\mathbf O$-module instead of left $\mathbf O$-module. The notions of a
homomorphism of $\mathbf O$-modules and of composition and addition of such
homomorphisms are evident. We thus obtain the \textit{additive category
$\mathbf C^{\mathbf O}$ of $\mathbf O$-modules on $X$}. For example, if $k$ is a ring
with identity, let $k_X$ be the corresponding constant sheaf of rings on
$X$. The category of $k_X$-modules is precisely the category of sheaves of
$k$-modules on $X$; for $k=\mathbf Z$, it is the category of sheaves of
abelian groups. Recall the following facts, already noted in 1.5 and 1.9.

\LayerTarget{tohoku:prop:3.1.1}
\textsc{Proposition 3.1.1.} \textit{Let $\mathbf O$ be a sheaf of rings
with identity on the space $X$. The additive category $\mathbf C^{\mathbf O}$ of
$\mathbf O$-modules on $X$ is an abelian category satisfying axioms \textup{AB 5)} and
\textup{AB 3*)}, and it has a generator.}

\LayerTarget{tohoku:proof:3.1.1:p1}
The direct sum $\mathrm S$ of a family $(F_i)$ of sheaves is constructed by taking,
for each open $U$, the direct sum of the $F_i(U)$ and then passing to the
sheaf associated to the resulting presheaf. The product $P$ of the sheaves
$F_i$ is constructed in the same way. The essential difference is that for
every $x\in X$, $\mathrm S(x)$ is the direct sum of the $F_i(x)$, whereas $P(x)$ \textit{is
not} the direct product of the $F_i(x)$. It is also easy to see that axiom
AB 4*) is not satisfied in general; for example, take $\mathbf O$ to be the
constant sheaf $\mathbf Z_X$. Finally, for every open $U\subset X$, let
$\mathbf O_U$ be the sheaf of $\mathbf O$-modules whose restriction to $\complement U$
is zero and whose restriction to $U$ agrees with $\mathbf O$; see [4, XVII,
prop. 1]. The family of the $\mathbf O_U$ is a system of generators of
$\mathbf C^{\mathbf O}$, as follows readily from Definition 1.9 and from the fact
that the homomorphisms from $\mathbf O_U$

\TohokuContinuation{156}

